% FailWhere — ICLR Workshop condensed version (4 pages, non-archival)
%
% Target: ICLR 2027 workshop on Evaluating and Aligning LLMs (~Feb 1, 2027)
% VERIFY target workshop is non-archival before submitting (to keep ARR submission clean)
%
% To compile: requires ICLR 2027 workshop template (typically iclr2027_workshop.sty).
% Or use standard article format for initial draft.

\documentclass[11pt]{article}
\usepackage{times}
\usepackage{latexsym}
\usepackage{microtype}
\usepackage{booktabs}
\usepackage{amsmath}
\usepackage{url}
\usepackage[margin=1in]{geometry}

\title{\textsc{FailWhere}: Verbal Failure Self-Localization in VLMs\\
is Generic Image Re-Reading with a Center Prior}

\author{Anonymous}

\begin{document}

\maketitle

\begin{abstract}
Can a Vision-Language Model identify which spatial region of an image caused
an incorrect answer?
We test this on the VSR dataset (142 failures, Claude Haiku, temp$=0$).
Four converging findings: the verbal \textsc{region:} tag
(1)~is produced equally well by a different model that never answered
(Sonnet-5-5, $+$25.4\,pp above null, $p{=}0.002$);
(2)~has equal accuracy on correct and failed answers ($\Delta{=}-1.1$\,pp,
TOST confirmed);
(3)~carries zero causal signal when center is tagged; and
(4)~lands $+$23--25\,pp above null when a specific non-center region is named.
The tag is generic image re-reading with a center prior, not failure self-knowledge.
Code and data: \url{https://anonymous.4open.science/r/vlm-failure-bench}.
\end{abstract}

\section{Introduction and Background}

Prior VLM attribution work is either white-box (attention, gradients) or focuses
on correct answers.
No prior work tests whether verbal self-reports can localize the spatial cause of
a \emph{failure}.
We introduce \textsc{FailWhere}: for each failure, ask the model to name the
responsible 3$\times$3 grid region using a structured \textsc{region:} tag.
Counterfactual masking (mask each of 9 cells with gray, re-run the model)
provides causal ground truth (CF-GT).

\section{Setup}

\textbf{Dataset.} VSR \cite{liu-2023-vsr}, 500 samples (142 failures), COCO images.
\textbf{Conditions.} Neutral (Cond.\,A, all items), Stateless (Cond.\,B, failures only),
Multi-turn (Cond.\,C, T/F exchange as context).
\textbf{CF-GT.} Per-failure masking: 51.4\% \texttt{none\_found}, 33.1\% \texttt{multiple},
15.5\% clean single-cell (hard attribution ceiling).
\textbf{Baselines.} Permuted-tag null (shuffle tags within-item); always-center;
COCO bbox hit (does the tag land on a known object bbox?).

\section{Results}

\begin{table}[h]
\centering
\small
\begin{tabular}{lrrrr}
\toprule
 & \multicolumn{2}{c}{Sonnet-5-5} & \multicolumn{2}{c}{Haiku-4-5} \\
 & Neutral & Stateless & Neutral & Stateless \\
\midrule
Clean hit ($n{=}22$)   & 31.8\% & \textbf{36.4\%} & 27.3\% & 31.8\% \\
Excess over null       & $+$20.4\,pp & $+$25.4\,pp & — & — \\
$p$ (one-sided)        & 0.007 & \textbf{0.002} & — & — \\
\midrule
\multicolumn{5}{l}{\emph{Two-regime split (Haiku stateless)}} \\
Tag = center ($n{=}3$)       & \multicolumn{4}{r}{excess $= 0.0$\,pp (all conditions)} \\
Tag $\neq$ center ($n{=}19$) & \multicolumn{4}{r}{excess $= +$16.4\,pp (all conditions)} \\
\midrule
\multicolumn{5}{l}{\emph{Correctness invariance (COCO bbox, $n{=}467$)}} \\
Subject hit (correct)  & \multicolumn{4}{r}{46.1\%} \\
Subject hit (failed)   & \multicolumn{4}{r}{45.0\% ($\Delta{=}-1.1$\,pp, TOST $p{=}0.86$)} \\
\bottomrule
\end{tabular}
\caption{Summary results. The \textsc{region:} tag does not distinguish failure
from success, and a different model matches the answering model's attribution
accuracy.}
\end{table}

\paragraph{Sonnet baseline (pillar~1).}
Sonnet-5-5 attributes Haiku's failures as well as Haiku does ($+$25.4\,pp, $p{=}0.002$).
No model-state access is required; the tag is generic image re-reading.

\paragraph{Correctness invariance (pillar~2).}
COCO bbox hit rates are within 2.6\,pp for correct vs.\ failed answers.
TOST equivalence ($\pm$8\,pp) confirms: the tag does not detect failure.

\paragraph{Two-regime structure (pillars 3--4).}
Center tags: 47\% of self-reports, 4.2$\times$ uniform prior, zero excess above null.
Non-center tags: $+$16--25\,pp above null, consistent across all three prompt
conditions and both models.
The tag is informative when the model escapes the center prior (${\sim}46\%$ of
cases), uninformative otherwise.

\section{Discussion and Conclusion}

Verbal failure self-localization in current VLMs is: prior-dominated (center),
not failure-specific (correctness-invariant), and not self-knowledge
(other-model control).
The positive finding is narrow but real: non-center self-reports are ${\sim}2\times$
above chance.
This motivates evaluation frameworks that score only non-center tags, and training
objectives that reward calibrated non-center commitment over center-defaulting.

\bibliographystyle{plain}
\bibliography{refs}

\end{document}
